\section{Preliminary Results}
\label{sec:research_method_preliminary_results}

Foundational progress has been achieved toward establishing the computational framework required for this investigation.  
A shared variational autoencoder (VAE) has been successfully trained on a combined dataset of normal and tuberculosis (TB) chest radiographs, providing a unified latent representation suitable for training latent diffusion models (LDMs).  
Using this shared encoder–decoder backbone, two unconditional LDMs were independently trained—one on normal chest X-rays and the other on TB X-rays—to serve as disease-specific generative priors.  
Initial compositional experiments conducted using the SuperDiff framework demonstrate that score-space superposition between these models can be meaningfully controlled: setting the weight of one model to zero reproduces the distribution of the other, effectively yielding a logical OR operation in generative space.  
These preliminary findings indicate that the models can generate plausible samples representative of their respective training domains, validating the soundness of the pipeline for subsequent compositional investigations.  
Figure~\ref{fig:progress_readiness_ldm} illustrates representative samples from the trained diffusion models, which serve as a readiness milestone for future experiments examining latent-space composition under both Euclidean and Riemannian diffusion formulations.

\begin{figure}[H]
    \centering
    % --- First image: Normal X-rays ---
    \begin{subfigure}[t]{0.48\textwidth}
        \centering
        \includegraphics[width=\textwidth]{chapters/research_method/media/normal_ldm_samples.png}
        \caption{Samples from the unconditional diffusion SDE trained on Normal chest X-rays.}
        \label{fig:normal_ldm_samples}
    \end{subfigure}
    \hfill
    % --- Second image: TB X-rays ---
    \begin{subfigure}[t]{0.48\textwidth}
        \centering
        \includegraphics[width=\textwidth]{chapters/research_method/media/tb_ldm_samples.png}
        \caption{Samples from the unconditional diffusion SDE trained on Tuberculosis (TB) chest X-rays.}
        \label{fig:tb_ldm_samples}
    \end{subfigure}

    % --- Unified caption ---
    \caption{
        Visual demonstration of the experimental readiness phase.
        Each panel shows images synthesized by separately trained unconditional diffusion SDEs, one on Normal and one on Tuberculosis (TB) chest X-ray datasets. 
        These models serve as proxies for studying compositionality in diffusion processes—specifically, whether latent manifolds of distinct disease distributions are geometrically and statistically composable under the SuperDiff framework.
        The ongoing investigation tests whether such manifolds admit coherent superposition through score-space operations, providing an empirical basis for evaluating the mathematical foundations of composable diffusion models.
    }
    \label{fig:progress_readiness_ldm}
\end{figure}